\section{Dynamical consequences}
\label{sec:consequences}

The spectral model gives ordinary mixing directly. The passage to higher
orders uses the separate multiple-mixing theorem; the entropy and Lyapunov
conclusions do not depend on that additional input.

\begin{corollary}\label{cor:dynamical-profile}
For the diffeomorphism $T$ and normalized standard volume $\mu$ of
Theorem~\ref{thm:main}, every integer $k\ge2$ and every fixed choice of
measurable sets $A_1,\ldots,A_k\subseteq\T^3$ satisfy
\[
 \mu\left(\bigcap_{i=1}^k T^{-t_i}A_i\right)
 \longrightarrow\prod_{i=1}^k\mu(A_i)
\]
whenever $t_1<\cdots<t_k$ are integers and
$\min_{1\le i<k}(t_{i+1}-t_i)\to\infty$. Moreover, the
Kolmogorov--Sinai entropy satisfies $h_\mu(T)=0$, and all three Lyapunov
exponents of $T$, counted with multiplicity, vanish at $\mu$-almost every
point.
\end{corollary}

\begin{proof}
Let $V:L^2_0(\T^3,\mu)\to L^2(S^1,m)$ be the unitary spectral model in
Theorem~\ref{thm:main}. For $a,b\in L^2_0(\T^3,\mu)$,
\[
 \int_{\T^3}(a\circ T^n)\overline b\dd\mu
 =\int_{S^1}w^n(Va)(w)\overline{(Vb)(w)}\dd m(w)
 \longrightarrow0\qquad(|n|\to\infty).
\]
Indeed, $(Va)\overline{Vb}\in L^1(m)$ by Cauchy--Schwarz, so the limit is
the Riemann--Lebesgue lemma. Taking centered indicators gives ordinary
mixing. The map $T$ is an invertible probability-preserving transformation
with measurable inverse, so the multiple-mixing theorem
\cite[Theorem~1.1]{MultipleMixing2026} gives every order $k\ge3$.
Invariance shifts $t_1$ to zero and identifies its positive time gaps with
$t_{i+1}-t_i$; the case $k=2$ is ordinary mixing.

The full Koopman operator is the direct sum of the one-dimensional
constants and the simple Lebesgue model on all of $L^2_0$, and hence has
finite spectral multiplicity. Normalized torus volume is a Lebesgue
probability space, and $T$ is a measure-preserving automorphism.
Rokhlin's finite-multiplicity entropy theorem
\cite[\S 7, p.~13]{Rokhlin1960} therefore gives $h_\mu(T)=0$.

The derivative and inverse derivative are bounded on the compact torus,
so the invertible Oseledets integrability conditions hold. Write
$\lambda_1(x),\lambda_2(x),\lambda_3(x)$ for the exponents repeated
according to multiplicity at regular points. Pesin's entropy formula,
as stated and reproved by Ma\~n\'e \cite[p.~95, equation~(2)]{Mane1981},
applies to the $C^\infty$ diffeomorphism and its absolutely continuous
invariant probability $\mu$. Thus
\[
 0=h_\mu(T)=\int_{\T^3}\sum_{i=1}^3\max\{\lambda_i(x),0\}\dd\mu(x).
\]
The nonnegative integrand vanishes almost everywhere, so every exponent
is nonpositive there. Preservation of standard volume gives
$|\det DT_x|=1$, and at every Oseledets-regular point
\[
 \sum_{i=1}^3\lambda_i(x)
 =\lim_{n\to\infty}\frac1n\log|\det D(T^n)_x|=0.
\]
The three nonpositive exponents must therefore all be zero at
$\mu$-almost every point.
\end{proof}
